% !PRINTED 366
\weberpage{366}
%<<CONT
that $x - x_0$ is divisible by $\alpha$; if accordingly we set, denoting by $h$ an arbitrary integer,
\[
x = x_0 + h \alpha,
\]
it follows that:
\[
y = y_0 + h \beta,
\]
and all solutions of (1) are contained in this form.

We therefore need only to seek \gesperrt{one} solution of (1), which we can always obtain in the following way. We assume $\beta$ to be positive, which does not essentially restrict the generality, since if need be we can change $x$ into $-x$, and, following \S.\,115, we develop the rational number $\alpha : \beta$ into a continued fraction:
\[
\frac{\alpha}{\beta} = (a_0, a_1, a_2 \ldots a_{n-1}).
\]

From this we form the convergents
\[
\frac{P_0}{Q_0}, \quad \frac{P_1}{Q_1}, \quad \frac{P_2}{Q_2} \cdots \frac{P_{n-1}}{Q_{n-1}}, \quad \frac{P_n}{Q_n},
\]
so that
\[
\frac{\alpha}{\beta} = \frac{P_n}{Q_n}.
\]

But since $\alpha$, $\beta$ as well as $P_n$, $Q_n$ are without common divisor, and since $\beta$ and $Q_n$ are positive, it follows that
\[
\alpha = P_n, \quad \beta = Q_n.
\]
Since further, by \S.\,117,
\[
P_n Q_{n-1} - Q_n P_{n-1} = (-1)^n
\]
it follows that

\weq{3}{x = (-1)^n P_{n-1}, \quad y = (-1)^n Q_{n-1}}
are integers which satisfy the equation (1), and with this the problem is completely solved.

For moderately large numbers the calculation is quite simple.

If, for example, we set $\alpha = 24335$, $\beta = 3588$, we first obtain by \S.\,115 the continued fraction:
\[
\frac{24335}{3588} = (6,\ 1,\ 3,\ 1,\ 1,\ 2,\ 6,\ 2,\ 1,\ 1,\ 4),
\]
and the convergents
\[
\frac{1}{0}, \quad \frac{6}{1}, \quad \frac{7}{1}, \quad \frac{27}{4}, \quad \frac{34}{5}, \quad \frac{61}{9}, \quad \frac{156}{23}, \quad \frac{997}{147}, \quad \frac{2150}{317}, \quad \frac{3147}{464},
\]
\[
\frac{5297}{781}, \quad \frac{24335}{3588},
\]
%CONT>>
